\appendixchapter{Б}{Программа и методика испытаний}

\begin{center}
\textbf{ПРОГРАММА И МЕТОДИКА ИСПЫТАНИЙ}\\
информационной системы\\
«Guide Helper»\\
\vspace{0.5em}
ЕСПД. ГОСТ 19.301-79~\cite{gost19301}
\end{center}

\section{Объект испытаний}

Объект испытаний~--- информационная система для создания и управления туристическими маршрутами с геопривязанными фотографиями «Guide Helper», версия 1.0.

Идентификационное наименование программы: \texttt{guide\_helper}.

Программа реализована в виде набора Docker-контейнеров и развёртывается средствами Docker Compose или Kubernetes.

\section{Цель испытаний}

Целью испытаний является проверка соответствия программы требованиям технического задания (Приложение~А), а именно:
\begin{itemize}
  \item проверка правильности реализации функциональных требований;
  \item проверка выполнения нефункциональных требований (производительность, надёжность, безопасность);
  \item выявление дефектов и несоответствий требованиям;
  \item подтверждение готовности программы к эксплуатации.
\end{itemize}

\section{Требования к программе}

Требования к функциональным характеристикам и надёжности приведены в техническом задании (Приложение~А, разделы~4.1 и~4.2).

\section{Требования к программной документации}

Для проведения испытаний необходимо наличие следующей документации:
\begin{itemize}
  \item настоящая программа и методика испытаний;
  \item техническое задание (Приложение~А);
  \item руководство пользователя (раздел~2.8 пояснительной записки).
\end{itemize}

\section{Средства и порядок испытаний}

\subsection{Средства испытаний}

Для проведения испытаний используются следующие средства:
\begin{itemize}
  \item сервер для развёртывания системы (требования~--- см. Приложение~А, раздел~4.4);
  \item клиентский компьютер с браузером Chrome~90+ или Firefox~89+;
  \item инструмент нагрузочного тестирования Apache JMeter;
  \item инструмент тестирования API Postman;
  \item встроенные средства тестирования языка Rust (\texttt{cargo test}).
\end{itemize}

\subsection{Порядок проведения испытаний}

Испытания проводятся в следующем порядке:
\begin{enumerate}
  \item Развёртывание системы в тестовой среде (\texttt{docker compose up -d}).
  \item Выполнение модульных тестов (\texttt{cargo test}).
  \item Выполнение функциональных тест-кейсов по методикам разделов~6.1--6.4.
  \item Проведение нагрузочного тестирования по методике раздела~6.5.
  \item Оформление протокола испытаний.
\end{enumerate}

\section{Методы испытаний}

\subsection{Тестирование аутентификации и авторизации}

\begin{spacing}{1.0}
\small
\begin{longtable}{|c|p{5.5cm}|p{3cm}|p{2.5cm}|}
\caption{Тест-кейсы аутентификации}
\label{tab:test-auth} \\
\hline
\textbf{№} & \textbf{Действие} & \textbf{Ожидаемый результат} & \textbf{Статус} \\
\hline
\endfirsthead
\multicolumn{4}{r}{\small\textit{Продолжение таблицы~\ref{tab:test-auth}}} \\
\hline
\endhead
А-01 & Регистрация с валидными данными (email, пароль $\geq$ 8~симв.) & Ответ 201, создана учётная запись & Выполнен \\
\hline
А-02 & Регистрация с существующим email & Ответ 409 Conflict & Выполнен \\
\hline
А-03 & Вход с корректными учётными данными & Ответ 200, получены access и refresh токены & Выполнен \\
\hline
А-04 & Вход с неверным паролем & Ответ 401 Unauthorized & Выполнен \\
\hline
А-05 & Обновление access-токена по refresh-токену & Ответ 200, новый access-токен & Выполнен \\
\hline
А-06 & Запрос к защищённому эндпоинту без токена & Ответ 401 Unauthorized & Выполнен \\
\hline
А-07 & Запрос с просроченным access-токеном & Ответ 401, сообщение об истечении токена & Выполнен \\
\hline
А-08 & Использование refresh-токена как access-токена & Ответ 401, неверный тип токена & Выполнен \\
\hline
\end{longtable}
\end{spacing}

\subsection{Тестирование управления маршрутами}

\begin{spacing}{1.0}
\small
\begin{longtable}{|c|p{5.5cm}|p{3cm}|p{2.5cm}|}
\caption{Тест-кейсы управления маршрутами}
\label{tab:test-routes} \\
\hline
\textbf{№} & \textbf{Действие} & \textbf{Ожидаемый результат} & \textbf{Статус} \\
\hline
\endfirsthead
\multicolumn{4}{r}{\small\textit{Продолжение таблицы~\ref{tab:test-routes}}} \\
\hline
\endhead
М-01 & Создание маршрута с 3 точками & Ответ 201, маршрут создан, возвращён UUID & Выполнен \\
\hline
М-02 & Получение маршрута по UUID & Ответ 200, полные данные маршрута & Выполнен \\
\hline
М-03 & Обновление маршрута (добавление точки) & Ответ 200, обновлённые данные & Выполнен \\
\hline
М-04 & Удаление маршрута владельцем & Ответ 204, маршрут удалён & Выполнен \\
\hline
М-05 & Попытка удаления маршрута другим пользователем & Ответ 403 Forbidden & Выполнен \\
\hline
М-06 & Публикация маршрута & Ответ 200, маршрут появился в каталоге & Выполнен \\
\hline
М-07 & Просмотр опубликованного маршрута по токену (без авторизации) & Ответ 200, данные маршрута & Выполнен \\
\hline
М-08 & Импорт маршрута из GeoJSON & Ответ 201, маршрут создан из файла & Выполнен \\
\hline
М-09 & Экспорт маршрута в GPX & Ответ 200, валидный GPX-файл & Выполнен \\
\hline
М-10 & Расчёт статистик (расстояние, высота, время) & Корректные значения по формулам & Выполнен \\
\hline
\end{longtable}
\end{spacing}

\subsection{Тестирование асинхронной обработки фотографий}

\begin{spacing}{1.0}
\small
\begin{longtable}{|c|p{5.5cm}|p{3cm}|p{2.5cm}|}
\caption{Тест-кейсы обработки фотографий}
\label{tab:test-photos} \\
\hline
\textbf{№} & \textbf{Действие} & \textbf{Ожидаемый результат} & \textbf{Статус} \\
\hline
\endfirsthead
\multicolumn{4}{r}{\small\textit{Продолжение таблицы~\ref{tab:test-photos}}} \\
\hline
\endhead
Ф-01 & Загрузка JPEG-фото (< 5~МБ) с GPS в EXIF & Статус «обрабатывается», уведомление по WebSocket & Выполнен \\
\hline
Ф-02 & Завершение обработки~--- оригинал в MinIO & URL оригинала в ответе API & Выполнен \\
\hline
Ф-03 & Завершение обработки~--- миниатюра 300px в MinIO & URL миниатюры в ответе API & Выполнен \\
\hline
Ф-04 & Время обработки файла 5~МБ & Менее 10~секунд & Выполнен \\
\hline
Ф-05 & Загрузка PNG-фото без EXIF & Обработка выполнена, координаты точки не изменены & Выполнен \\
\hline
Ф-06 & Загрузка файла неподдерживаемого формата & Ответ 400 Bad Request & Выполнен \\
\hline
Ф-07 & Загрузка файла > 10~МБ & Ответ 413 Payload Too Large & Выполнен \\
\hline
\end{longtable}
\end{spacing}

\subsection{Тестирование ИИ-ассистента}

\begin{spacing}{1.0}
\small
\begin{longtable}{|c|p{5.5cm}|p{3cm}|p{2.5cm}|}
\caption{Тест-кейсы ИИ-ассистента}
\label{tab:test-ai} \\
\hline
\textbf{№} & \textbf{Действие} & \textbf{Ожидаемый результат} & \textbf{Статус} \\
\hline
\endfirsthead
\multicolumn{4}{r}{\small\textit{Продолжение таблицы~\ref{tab:test-ai}}} \\
\hline
\endhead
И-01 & Отправка текстового сообщения ассистенту & Получен ответ, создан диалог & Выполнен \\
\hline
И-02 & Запрос геокодирования («найди Эльбрус») & Ассистент вызвал инструмент, вернул координаты & Выполнен \\
\hline
И-03 & Запрос поиска маршрутов («покажи лёгкие маршруты») & Ассистент вернул список маршрутов & Выполнен \\
\hline
И-04 & Запрос погоды для точки маршрута & Ассистент вернул прогноз погоды & Выполнен \\
\hline
И-05 & Превышение лимита сообщений (rate limit) & Ответ 429 Too Many Requests & Выполнен \\
\hline
\end{longtable}
\end{spacing}

\subsection{Нагрузочное тестирование}

Нагрузочное тестирование проводится с использованием Apache JMeter. Параметры теста:

\begin{itemize}
  \item количество виртуальных пользователей: 100;
  \item период нарастания нагрузки: 60~секунд;
  \item продолжительность теста: 10~минут;
  \item целевые эндпоинты: \texttt{GET /api/v1/routes} (каталог маршрутов), \texttt{GET /api/v1/routes/\{id\}} (данные маршрута).
\end{itemize}

Критерии приёмки нагрузочного тестирования:

\begin{itemize}
  \item среднее время отклика: не более 200~мс;
  \item 95-й перцентиль времени отклика: не более 500~мс;
  \item процент ошибок: не более 1\%;
  \item пропускная способность: не менее 100~запросов/сек.
\end{itemize}

По результатам проведённого нагрузочного тестирования: среднее время отклика составило 80~мс, 95-й перцентиль~--- 210~мс, процент ошибок~--- 0\%, пропускная способность~--- 145~запросов/сек. Все критерии выполнены.
